\section{Small corollaries for \texorpdfstring{$k=7$}{k=7} and \texorpdfstring{$k=8$}{k=8}}\label{ec:sec}
% Local to this section: the statement heads contain long unbreakable Lean names; \sloppy lets TeX break them at spaces.
\begingroup\sloppy

This section gives written proofs for a group of short Lean files about rotationally symmetric point sets.
\begin{enumerate}[label=(\arabic*)]
\item \emph{Conditional $ES(7)$ corollaries} (Subsection~\ref{ec:ssec:es7}). If the case classes
  $\mathrm{Fivefold35}$ and $\mathrm{Fourfold32}$ hold (they are hypotheses of the Lean theorems), then every finite
  set of points in general position that is invariant under the rotation by $2\pi/5$ about some point and has at
  least $32$ points, and every such set for the rotation by $2\pi/4$ with at least $29$ points, has $7$ points in
  convex position.
\item \emph{Rotations of order at least $8$} (Subsection~\ref{ec:ssec:order8}). A subset of $\C$ that is mapped into
  itself by $z\mapsto e^{2\pi i/m}z$ with $m\ge8$ and contains no $8$ points in convex position is contained in
  $\{0\}$. Also the threshold arithmetic for $m=4,5,6,7$: for all $t,e\in\N$ with $e\le1$, $65\le4t\iff t\ge17$,
  $65\le5t+e\iff t\ge13$, $65\le6t\iff t\ge11$ and $65\le7t+e\iff t\ge10$.
\item \emph{Orientations in inverse-radius polar form} (Subsection~\ref{ec:ssec:polar}): the sign of the orientation of
  three points $(\cos\varphi_k,\sin\varphi_k)/u_k$, and its consequences for layered configurations and for nested
  regular $m$-gons.
\item \emph{Size compatibility} (Subsections~\ref{ec:ssec:compat} and~\ref{ec:ssec:l0}): the orbit count
  $\card S\equiv[0\in S]\pmod m$ for finite sets invariant under an element of exact order $m$, its converse under
  explicit hypotheses, the classification of the $m$ with $2^{m-1}\bmod m\in\{0,1\}$, and the consequence for
  $m$-fold symmetric sets of $2^{m-1}$ points.
\end{enumerate}

\paragraph{Sources.} Paths are relative to \dirn{cloud/es7-2026-10-19/results}. We computed the SHA-256 of each file;
the table gives the first $12$ hexadecimal digits. The line ranges are the scope of this section; the Coverage paragraph
at the end lists what in them is not covered.
\begin{center}\footnotesize
\begin{tabular}{@{}lll@{}}
\toprule
Lean file (lines in scope) & SHA-256 & axioms log\\
\midrule
\dirn{LEANEXT/es7corr/ES7Corollaries.lean} (1148--1233) & \texttt{5883c8ebac55} & \dirn{LEANEXT/es7corr/axioms.log}\\
\dirn{ES8SYM/lean/C6check.lean} (1--123) & \texttt{92f15b5cc907} & \dirn{ES8SYM/lean/C6check.out}\\
\dirn{ES8SYM/lean/PolarOrient.lean} (1--132) & \texttt{6e54f5f83a19} & \dirn{ES8SYM/lean/PolarOrient.out}\\
\dirn{PRIMEFOLD/compat/lean/PCompatCore.lean} (1--291) & \texttt{ebc6aacefadb} & \dirn{PRIMEFOLD/compat/lean/PCompatCore.out}\\
\dirn{PRIMEFOLD/compat/lean/L0Check.lean} (1--152) & \texttt{05a67c3b8b16} & \dirn{PRIMEFOLD/compat/lean/L0Check.out}\\
\bottomrule
\end{tabular}
\end{center}
According to the records, all five files were compiled with the repository's pins, Lean~4.24.0 and Mathlib
\texttt{f897ebcf}; nothing was re-run for this document. The hashes of \dirn{LEANEXT/es7corr/ES7Corollaries.lean} and
of the two Lean files in \dirn{PRIMEFOLD/compat/lean} agree with the prefixes recorded in \dirn{LEANEXT/es7corr/REPORT.md}
and \dirn{PRIMEFOLD/README.md}. For the two \dirn{ES8SYM} files, \dirn{ES8SYM/README.md} (lines 258--260) records the
prefixes \texttt{8c090168} and \texttt{4d42e2fb} for the compiled files and states that the present copies differ from
them only in four comment lines; we could not check this, so the logs of these two files refer to the compiled
versions.

\paragraph{Axioms.} The logs list the following theorems, each depending on at most the standard axioms
\lean{propext}, \lean{Classical.choice} and \lean{Quot.sound}: the five $ES(7)$ corollaries,
Theorems~\ref{ec:thm:fivefold-erase}--\ref{ec:thm:fourfold} (\dirn{LEANEXT/es7corr/axioms.log}); \lean{RefuteC6.part_a}
and \lean{RefuteC6.thr_m4}--\lean{thr_m7} (\dirn{ES8SYM/lean/C6check.out}); the six theorems \lean{C7.orient_polar},
\lean{orient_polar_cross}, \lean{orient_layered}, \lean{orient_layered_scale}, \lean{nested_regular_orient},
\lean{orient_centre} (\dirn{ES8SYM/lean/PolarOrient.out}); \lean{PCompat.dvd_card}, \lean{card_mod},
\lean{exists_iff}, \lean{compat_even}, \lean{compat_odd}, \lean{classification}, and \lean{PCompat.ex341}, which
depends on \lean{propext} only (\dirn{PRIMEFOLD/compat/lean/PCompatCore.out}); \lean{L0Check.L0_rot}
(\dirn{PRIMEFOLD/compat/lean/L0Check.out}). Every other declaration covered here occurs in the statement or the proof
of one of these, so the same logs cover it, except \lean{RefuteC6.m4_with_centre_counter}, which has no axiom line (it is a closed
arithmetic statement whose Lean proof is \lean{omega}; the compile log of its file shows no error). No file uses \lean{sorry},
\lean{native_decide} or new axioms; \lean{PCompat.ex341} is checked by kernel evaluation (\lean{decide}).

\paragraph{Rigor and dependencies.} The theorems with a Lean reference are Lean-checked, with the axioms above; for the
two \dirn{ES8SYM} files this rests on the logs of the compiled versions, which according to the records differ from the
present copies only in comment lines (see Sources). The proofs below are written proofs (not refereed). Statements
tagged ``written step'' are not declarations of the five Lean files: Fact~\ref{ec:fact:prim} restates Mathlib lemmas,
and Corollary~\ref{ec:cor:planar} has no Lean counterpart. The theorems of Subsection~\ref{ec:ssec:es7} are
conditional on the propositions $\mathrm{Fivefold35}$ and $\mathrm{Fourfold32}$, and Lemma~\ref{ec:lem:pc-exist} and
Theorem~\ref{ec:thm:compat}(c) take hypotheses on $r$ that Lean does not discharge (Remark~\ref{ec:rem:compat}). We
use Theorem~\ref{ll:thm:reduction} (reduction to exactly $mt$ points) and Lemma~\ref{lg:lem:centre} (even orders
exclude the centre), and the conventions of Definitions~\ref{fd:def:point}, \ref{fd:def:orient}, \ref{fd:def:gp},
\ref{fd:def:cp}, \ref{fd:def:ngon}, \ref{fd:def:inv} and~\ref{fd:def:rot}; remarks also refer to
Lemma~\ref{ll:lem:minperiod}, Theorems~\ref{ll:thm:layers} and~\ref{ll:thm:card-peel}, Lemma~\ref{lg:lem:polar},
Lemmas~\ref{oc:lem:rot-inv}, \ref{oc:lem:perm} and~\ref{oc:lem:same-ray}, Theorem~\ref{oc:thm:phase-orbit} and
Definition~\ref{sr:def:classes}. The other subsections are self-contained apart from standard facts, which we state
where they are used.

%---------------------------------------------------------------------------------------------------------------------
\subsection{Conditional \texorpdfstring{$ES(7)$}{ES(7)} corollaries for fivefold and fourfold symmetry}\label{ec:ssec:es7}

The source is Section~11 of \dirn{LEANEXT/es7corr/ES7Corollaries.lean} (lines 1148--1233, namespace
\lean{ErdosSzekeres.LayerLemma}). Lines 1--1147 of this file coincide line by line with lines 1--1147 of
\dirn{LAYERLEAN/LayerLemma.lean} (we checked this with \texttt{diff}; the latter file has one more line, its closing
\texttt{end}). Hence the declarations \lean{Invariant}, \lean{rotAbout}, \lean{centre_not_mem_of_even} and
\lean{convexNGon_of_layers} used below are verbatim copies of those of \dirn{LAYERLEAN/LayerLemma.lean}, at the same
lines ($105$, $298$, $582$, $1066$).

\emph{Conventions.} $S$ is a finite set of points (Lean \texttt{Finset Point}) and $c$ a point. $\mathrm{rotAbout}(c,\theta)$
is the counterclockwise rotation by $\theta$ about $c$ (Definition~\ref{fd:def:rot}). ``$S$ is invariant under $f$''
means $f(p)\in S$ for every $p\in S$ (Definition~\ref{fd:def:inv}). General position is
Definition~\ref{fd:def:gp}, and $\mathrm{HasConvexNGon}(n,S)$ means that some $n$-element subset of $S$ is in convex
position (Definition~\ref{fd:def:ngon}; no emptiness condition). We write $S\setminus\{c\}$ for Lean's
\texttt{S.erase c}. In the statements below the angles are $2\pi/5$ and $2\pi/4$ with real numerals $5$ and $4$; in
Theorem~\ref{ll:thm:reduction} and Lemma~\ref{lg:lem:centre} (items (R) and (C) below) the angle is $2\pi/m$ with
$m\in\N$ cast to $\R$. For $m=5$ and $m=4$ these are the same real numbers (the Lean proofs pass between the two forms
with the cast-normalizing tactic \lean{exact_mod_cast}).

\begin{definition}[the case classes; \leanref{ErdosSzekeres.LayerLemma.Fivefold35}{LEANEXT/es7corr/ES7Corollaries.lean:1157}, \leanref{ErdosSzekeres.LayerLemma.Fourfold32}{LEANEXT/es7corr/ES7Corollaries.lean:1164}]\label{ec:def:cases}
$\mathrm{Fivefold35}$ is the proposition: for every finite set $S$ of points and every point $c$, if $S$ is in
general position, $\card S=35$, $c\notin S$, and $S$ is invariant under $\mathrm{rotAbout}(c,2\pi/5)$, then
$\mathrm{HasConvexNGon}(7,S)$.

$\mathrm{Fourfold32}$ is the proposition: for every finite set $S$ of points and every point $c$, if $S$ is in general
position, $\card S=32$, $c\notin S$, and $S$ is invariant under $\mathrm{rotAbout}(c,2\pi/4)$, then
$\mathrm{HasConvexNGon}(7,S)$.
\end{definition}

\begin{remark}[status of the case classes]\label{ec:rem:cases}
No Lean theorem proves $\mathrm{Fivefold35}$ or $\mathrm{Fourfold32}$. The theorems of this subsection take them as
hypotheses and are therefore conditional. In the project they are supported by computer-assisted results outside
Lean, which also rest on written arguments. For SAT certificates the project counts only the verdicts of the proof
checker cake\_lpr (\dirn{WRITEUP/SYMMETRIC.md}, Section~3). The records also show acceptances by a second checker,
lrat-check, which the project does not count: it is a secondary check only, and its unpatched upstream build is
unsound on the empty clause (\dirn{E2/logs/lrat_check_soundness.md}; \dirn{LRATCHECK20261105}).
\begin{itemize}
\item $\mathrm{Fivefold35}$: \dirn{SYM5} (\dirn{SYM5/README.md}). Two formulas for the $35$-point fivefold case, an
  angle form and a layer form of one SAT encoder (not independent encodings: audit \dirn{AUDIT3}, finding A4-3), are
  recorded as unsatisfiable; the LRAT proofs of both were accepted by cake\_lpr. The passage from
  these verdicts to $\mathrm{Fivefold35}$ uses the soundness of the encoder (recorded as a written proof, reviewed in
  the audit, \dirn{AUDIT3}, A4), a normal-form lemma for the labelling (for the layer form, the layer lemma), and the
  centre-perturbation lemma, which makes the centre generic (\dirn{C2ENC/perturbation.md}, recorded there as a
  complete written proof, not Lean-checked, refereed within the project on the day it was written:
  \dirn{C2ENC/referee/NOTES.md}, Section~4). \dirn{C2ENC/class_table.md} summarizes this chain for the constant of
  the same text in \dirn{SYMLEAN} (see below).
\item $\mathrm{Fourfold32}$: \dirn{C4CUT} (\dirn{C4CUT/README.md}). The layered formula for the $32$-point fourfold
  case with six realizability cuts, each justified by a Grassmann--Pl\"ucker certificate checked in exact rational
  arithmetic, is recorded as unsatisfiable; its LRAT proof was accepted by cake\_lpr. The final run does not record
  which checker builds it used (\dirn{AUDIT3}, A4-5, which concludes that the claim is carried by the cake\_lpr verdict);
  the replay \dirn{REPLAY}, with hash-checked files and checker binaries, accepted the proof again with cake\_lpr. The
  passage to $\mathrm{Fourfold32}$ uses the soundness of the layered encoder (recorded in \dirn{C4CUT/README.md} as a
  written proof, refereed within the project: the review of 1~November 2026, \dirn{C4CUT/review/README.md}), the
  layer lemma, and the sector normal form: each of the other orbits has exactly one
  point in the open quadrant between the rays from the centre through two consecutive points of the outermost orbit.
  The audit found the argument for this not written down (\dirn{AUDIT3}, A4-1); \dirn{C4CUT/README.md}, Section~1,
  now gives it in one sentence: a point of another orbit on a bounding ray would be collinear with two antipodal
  points of the outermost orbit.
\end{itemize}
The reviews, audits and refereeing named in this list were all made within the project (internal reviews); none of
them is an external peer review. Apart from the layer lemma and the ray fact for even orders behind the sector
normal form, none of these written arguments is checked in Lean. Both have Lean-checked forms. For the layer lemma
it is Theorem~\ref{ll:thm:layers}.
For the ray fact it is Lemma~\ref{oc:lem:same-ray} together with Theorem~\ref{oc:thm:phase-orbit}
(\dirn{LEANEXT/orbit/OrbitConsistency.lean}): for even $m>0$, no two points of a finite set in general position that
is invariant under $\mathrm{rotAbout}(c,2\pi/m)$ lie on one open ray from $c$; and if $x$ and $y$ are points of such a
set, both different from $c$, and $y$ is not in the orbit of $x$ under that rotation, then the directions of $x$ and
$y$ at $c$ differ by no integer multiple of $2\pi/m$. The
remaining step, that each of the other orbits has exactly one point in the open quadrant, is a counting argument and
is not in Lean. The rotation invariance of orientations and their sign rule under permutations, on which the
encoders' single variable per orbit of triples rests, are Lean-checked too (Lemmas~\ref{oc:lem:rot-inv}
and~\ref{oc:lem:perm}). Nothing in Lean connects any of these Lean-checked forms to the formulas, and no Lean theorem
links a formula or a certificate to these propositions. Propositions with the same text occur in
\dirn{SYMLEAN/SymmetricReduction.lean:477} and \dirn{SYMLEAN/SymmetricReduction.lean:496}
(Definition~\ref{sr:def:classes}); they are different Lean constants (namespace \lean{ErdosSzekeres.SymReduction},
built on that file's own copies of \lean{Invariant} and \lean{rotAbout}), and no Lean theorem identifies the two
versions. The docstring of \lean{Fourfold32}, copied from \dirn{SYMLEAN}, says that it is used only by a variant
theorem of \dirn{SYMLEAN}; in \dirn{LEANEXT/es7corr/ES7Corollaries.lean} it is the hypothesis of
Theorems~\ref{ec:thm:fourfold-erase} and~\ref{ec:thm:fourfold}.
\end{remark}

We use two Lean-checked results whose written proofs are in other sections, in the following form.
\begin{itemize}
\item[(R)] Theorem~\ref{ll:thm:reduction} (\lean{convexNGon_of_layers}, \dirn{LAYERLEAN/LayerLemma.lean:1066}). Let
  $m,K,t\in\N$ with $0<m$ and $K\le2m$, and assume: \emph{for every finite set $T$ and every point $c'$, if $T$ is in
  general position, $c'\notin T$, $T$ is invariant under $\mathrm{rotAbout}(c',2\pi/m)$ and $\card T=m\cdot t$, then
  $\mathrm{HasConvexNGon}(K,T)$}. Then for every finite set $S$ and every point $c$ such that $S$ is in general
  position, $S$ is invariant under $\mathrm{rotAbout}(c,2\pi/m)$ and $m\cdot(t\mathbin{\dot{-}}1)<\card{S\setminus\{c\}}$,
  we have $\mathrm{HasConvexNGon}(K,S)$. Here $t\mathbin{\dot{-}}1=\max(t-1,0)$ is Lean's truncated subtraction
  on $\N$ (the notation of the layer-lemma section).
\item[(C)] Lemma~\ref{lg:lem:centre} (\lean{centre_not_mem_of_even}, \dirn{LAYERLEAN/LayerLemma.lean:582}). Let
  $m\in\N$, let $S$ be a finite set of points and $c$ a point. If $0<m$, $m$ is even, $S$ is invariant under
  $\mathrm{rotAbout}(c,2\pi/m)$ and in general position, and some $p\in S$ satisfies $p\ne c$, then $c\notin S$. (In
  this section, for a map $f$ and $n\in\N$, $f^n$ is the $n$-fold composite of $f$, written $f^{[n]}$ in the
  foundations section; so $f^0=\mathrm{id}$. The power $\mathrm{rotAbout}(c,2\pi/m)^{m/2}$ is the half-turn
  $p\mapsto 2c-p$, and $p$, $c$, $2c-p$ would be three distinct collinear points of $S$.)
\end{itemize}

In each of the following five theorems the case class is a hypothesis (Definition~\ref{ec:def:cases} and
Remark~\ref{ec:rem:cases}).

\begin{theorem}[fivefold, layer form; \leanref{ErdosSzekeres.LayerLemma.fivefold_heptagon_of_erase}{LEANEXT/es7corr/ES7Corollaries.lean:1170}]\label{ec:thm:fivefold-erase}
Assume $\mathrm{Fivefold35}$. Let $S$ be a finite set of points and $c$ a point such that $S$ is in general position,
$S$ is invariant under $\mathrm{rotAbout}(c,2\pi/5)$, and $30<\card{S\setminus\{c\}}$. Then
$\mathrm{HasConvexNGon}(7,S)$.
\end{theorem}

\begin{proof}
We follow the Lean proof. Apply (R) with $m=5$, $K=7$, $t=7$; here $0<5$ and $K=7\le10=2m$. The assumption of (R)
holds: let $T$ be finite and $c'$ a point such that $T$ is in general position, $c'\notin T$, $T$ is invariant under
$\mathrm{rotAbout}(c',2\pi/5)$ and $\card T=5\cdot7=35$. These are the four hypotheses of $\mathrm{Fivefold35}$ for
$T$ and $c'$ (listed there in a different order), so $\mathrm{Fivefold35}$ gives $\mathrm{HasConvexNGon}(7,T)$. The
size condition of (R) is $5\cdot(7\mathbin{\dot{-}}1)=5\cdot6=30<\card{S\setminus\{c\}}$, which is our hypothesis. So (R) gives
$\mathrm{HasConvexNGon}(7,S)$.
\end{proof}

\begin{theorem}[fivefold; \leanref{ErdosSzekeres.LayerLemma.fivefold_heptagon}{LEANEXT/es7corr/ES7Corollaries.lean:1183}]\label{ec:thm:fivefold}
Assume $\mathrm{Fivefold35}$. Let $S$ be a finite set of points and $c$ a point such that $S$ is in general position,
$S$ is invariant under $\mathrm{rotAbout}(c,2\pi/5)$, and $32\le\card S$. Then $\mathrm{HasConvexNGon}(7,S)$.
\end{theorem}

\begin{proof}
As in Lean: removing one element lowers the cardinality by at most one, so
$\card{S\setminus\{c\}}\ge\card S-1\ge31>30$, and Theorem~\ref{ec:thm:fivefold-erase} applies.
\end{proof}

\begin{theorem}[centre-free; \leanref{ErdosSzekeres.LayerLemma.fivefold_heptagon_of_centre_not_mem}{LEANEXT/es7corr/ES7Corollaries.lean:1193}]\label{ec:thm:fivefold-nc}
Assume $\mathrm{Fivefold35}$. Let $S$ be a finite set of points and $c$ a point such that $S$ is in general position,
$S$ is invariant under $\mathrm{rotAbout}(c,2\pi/5)$, $c\notin S$, and $31\le\card S$. Then
$\mathrm{HasConvexNGon}(7,S)$.
\end{theorem}

\begin{proof}
As in Lean: since $c\notin S$, $S\setminus\{c\}=S$, so $\card{S\setminus\{c\}}=\card S\ge31>30$, and
Theorem~\ref{ec:thm:fivefold-erase} applies.
\end{proof}

\begin{theorem}[fourfold, layer form; \leanref{ErdosSzekeres.LayerLemma.fourfold_heptagon_of_erase}{LEANEXT/es7corr/ES7Corollaries.lean:1203}]\label{ec:thm:fourfold-erase}
Assume $\mathrm{Fourfold32}$. Let $S$ be a finite set of points and $c$ a point such that $S$ is in general position,
$S$ is invariant under $\mathrm{rotAbout}(c,2\pi/4)$, and $28<\card{S\setminus\{c\}}$. Then
$\mathrm{HasConvexNGon}(7,S)$.
\end{theorem}

\begin{proof}
We follow the Lean proof. Apply (R) with $m=4$, $K=7$, $t=8$; here $0<4$ and $K=7\le8=2m$. The assumption of (R)
holds by $\mathrm{Fourfold32}$, exactly as in the proof of Theorem~\ref{ec:thm:fivefold-erase}, since $4\cdot8=32$.
The size condition of (R) is $4\cdot(8\mathbin{\dot{-}}1)=4\cdot7=28<\card{S\setminus\{c\}}$, our hypothesis.
\end{proof}

\begin{theorem}[fourfold; \leanref{ErdosSzekeres.LayerLemma.fourfold_heptagon}{LEANEXT/es7corr/ES7Corollaries.lean:1216}]\label{ec:thm:fourfold}
Assume $\mathrm{Fourfold32}$. Let $S$ be a finite set of points and $c$ a point such that $S$ is in general position,
$S$ is invariant under $\mathrm{rotAbout}(c,2\pi/4)$, and $29\le\card S$. Then $\mathrm{HasConvexNGon}(7,S)$.
\end{theorem}

\begin{proof}
We follow the Lean proof. Since $\card S\ge29>1$, $S$ is not contained in $\{c\}$, so some $p\in S$ satisfies $p\ne
c$. Apply (C) with $m=4$: $0<4$, $4$ is even, $S$ is in general position, and $S$ is invariant under
$\mathrm{rotAbout}(c,2\pi/4)$, which is the rotation $\mathrm{rotAbout}(c,2\pi/m)$ of (C) for $m=4$ (the cast angle
$2\pi/m$ equals $2\pi/4$ for $m=4$; see the Conventions). So $c\notin S$. Hence $S\setminus\{c\}=S$ and
$\card{S\setminus\{c\}}=\card S\ge29>28$, and
Theorem~\ref{ec:thm:fourfold-erase} applies.
\end{proof}

\begin{remark}[what the theorems give]\label{ec:rem:es7-sharp}
Under $\mathrm{Fivefold35}$, Theorem~\ref{ec:thm:fivefold} says that a finite set in general position that is
invariant under $\mathrm{rotAbout}(c,2\pi/5)$ for some $c$ and has no $7$ points in convex position has at most $31$
points; under $\mathrm{Fourfold32}$, Theorem~\ref{ec:thm:fourfold} gives at most $28$ points for
$\mathrm{rotAbout}(c,2\pi/4)$. These are the upper bounds $N(C_5)\le31$ and $N(C_4)\le28$ of the symmetric $ES(7)$
note (\dirn{WRITEUP/note/symmetric_es7.tex}, the proposition labelled \texttt{prop:C}), where they are derived from
the $35$-point fivefold and the $32$-point fourfold case by orbit counting (the lemma labelled \texttt{lem:sizes}). A set
with symmetry $C_s$ in the sense of the note is invariant under a rotation $\rho$ of order $s$ about some point $c$,
that is $\rho=\mathrm{rotAbout}(c,2\pi k/s)$ for an integer $k$ with $\gcd(k,s)=1$. Choose $j\in\N$ with
$jk\equiv1\pmod s$. By the complex form of Definition~\ref{fd:def:rot},
$\rho^j=\mathrm{rotAbout}(c,2\pi jk/s)=\mathrm{rotAbout}(c,2\pi/s)$, since the two angles differ by a multiple of
$2\pi$; and a set mapped into itself by $\rho$ is mapped into itself by $\rho^j$. So the two theorems apply to every
finite set in general position with symmetry $C_5$, respectively $C_4$. The matching lower-bound examples \dirn{SYMMAX/c5_31points.txt} and
\dirn{SYMMAX/c4_28points.txt} rest on exact computations outside Lean, which were not repeated for this document. The
logs \dirn{LEANEXT/es7corr/mutant31.log} and \dirn{LEANEXT/es7corr/mutant28.log} record that the Lean proofs fail at their
final arithmetic step when the thresholds $32$ and $29$ are lowered to $31$ and $28$; this audits the proof scripts
only.
\end{remark}

\begin{remark}[a route without the layer lemma]\label{ec:rem:orbits}
Theorems~\ref{ec:thm:fivefold-erase} and~\ref{ec:thm:fourfold-erase} also follow by orbit counting, the argument
of the symmetric $ES(7)$ note. Let $\rho=\mathrm{rotAbout}(c,2\pi/m)$ with $m\in\{4,5\}$. The map $\rho$ fixes $c$ and
is injective, so it maps $S'=S\setminus\{c\}$ into itself, hence onto itself ($S'$ is finite). Every $y\ne c$ has
minimal period $m$ under $\rho$ (Lemma~\ref{ll:lem:minperiod}), so the $\rho$-orbits $\{y,\rho y,\dots,\rho^{m-1}y\}$
of the points $y\in S'$ have exactly $m$ elements and partition $S'$. If $m(t-1)<\card{S'}$, with $t=7$ for $m=5$ and
$t=8$ for $m=4$, then $S'$ consists of at least $t$ orbits; the union $T$ of $t$ of them has $mt$ points ($35$,
respectively $32$), satisfies $c\notin T$, is in general position (as a subset of $S$) and is invariant under
$\rho$. The case class applies to $T$, and $7$ points of $T$ in convex position lie in $S$. The Lean proofs use
Theorem~\ref{ll:thm:reduction} instead, which takes for $T$ the $t$ outermost convex layers of $S'$.
\end{remark}

%---------------------------------------------------------------------------------------------------------------------
\subsection{Rotations of order at least eight}\label{ec:ssec:order8}

The source is \dirn{ES8SYM/lean/C6check.lean} (namespace \lean{RefuteC6}; it imports Mathlib only, not the project
library). It is the Lean check of a complex form of part (a) of Lemma~2 of \dirn{ES8SYM/proofs/layers.md} (the planar
form is Corollary~\ref{ec:cor:planar}, a written step) and of the arithmetic of part (b) of that lemma.

\emph{Conventions.} In this subsection points are complex numbers. $\C$ is a real vector space of dimension $2$,
$\abs z$ is the modulus (the Euclidean norm), and for $r\in\R$ we write $B_r=\{z\in\C:\abs z\le r\}$ (Lean
\texttt{closedBall 0 r}); the circle $\{z\in\C:\abs z=r\}$ is Lean's \texttt{sphere 0 r}, and both sets are empty if
$r<0$. $\conv$ is the convex hull over $\R$, the smallest convex set containing the given set. The following notions are
Mathlib's.
\begin{itemize}
\item An \emph{extreme point} of $A\subseteq\C$ (Mathlib \lean{Set.extremePoints}) is a point $x\in A$ such that
  whenever $x_1,x_2\in A$, $a,b\in\R$, $a>0$, $b>0$, $a+b=1$ and $x=ax_1+bx_2$, we have $x_1=x$.
\item A family $f:\iota\to\C$ is \emph{convex independent} over $\R$ (Mathlib \lean{ConvexIndependent}) if for every
  set $s\subseteq\iota$ and every $x\in\iota$, $f(x)\in\conv f(s)$ implies $x\in s$.
\item $\zeta\in\C$ is a \emph{primitive $m$-th root of unity} (Mathlib \lean{IsPrimitiveRoot}) if $\zeta^m=1$ and,
  for every $l\in\N$, $\zeta^l=1$ implies $m\mid l$.
\end{itemize}
$\mathrm{Fin}\,8=\{0,1,\dots,7\}$. A finite set of points is in convex position (Definition~\ref{fd:def:cp}) if the
inclusion family of the set is convex independent; Corollary~\ref{ec:cor:planar} makes the passage between injective
families in $\C$ and finite sets of points precise.

\begin{lemma}[points of a circle are extreme points of the disc; \leanref{RefuteC6.sphere_subset_extremePoints}{ES8SYM/lean/C6check.lean:24}]\label{ec:lem:sphere}
For every $r\in\R$, every $x\in\C$ with $\abs x=r$ is an extreme point of $B_r$.
\end{lemma}

\begin{proof}
We follow the Lean proof, which uses the strict convexity of the norm of $\C$; we make it explicit. For
$x_1,x_2\in\C$ and real $a,b$ with $a+b=1$,
\begin{equation*}
\abs{ax_1+bx_2}^2=a\abs{x_1}^2+b\abs{x_2}^2-ab\,\abs{x_1-x_2}^2 .
\end{equation*}
Indeed, with $\langle z,w\rangle=\operatorname{Re}(z\bar w)$ the left side is
$a^2\abs{x_1}^2+b^2\abs{x_2}^2+2ab\langle x_1,x_2\rangle$, the right side is
$(a-ab)\abs{x_1}^2+(b-ab)\abs{x_2}^2+2ab\langle x_1,x_2\rangle$, and $a-ab=a(1-b)=a^2$, $b-ab=b^2$.

Let $\abs x=r$. Then $x\in B_r$. Let $x_1,x_2\in B_r$ and $a,b>0$ with $a+b=1$ and $x=ax_1+bx_2$. If $x_1=x_2$, then
$x=(a+b)x_1=x_1$. Suppose $x_1\ne x_2$. Since $0\le\abs{x_1}\le r$ we have $r\ge0$ and $\abs{x_1}^2\le r^2$, and
likewise $\abs{x_2}^2\le r^2$; also $ab\abs{x_1-x_2}^2>0$. The identity gives
$r^2=\abs x^2<ar^2+br^2=r^2$, a contradiction. Hence $x_1=x$.
\end{proof}

\begin{lemma}[injective families on a circle are convex independent; \leanref{RefuteC6.convexIndependent_of_norm_eq}{ES8SYM/lean/C6check.lean:43}]\label{ec:lem:ci-circle}
Let $\iota$ be any type, $f:\iota\to\C$ injective, and $r\in\R$ with $\abs{f(i)}=r$ for every $i\in\iota$. Then $f$ is
convex independent over $\R$.
\end{lemma}

\begin{proof}
Step 1. \emph{If $x$ is an extreme point of a convex set $A\subseteq\C$, then $A\setminus\{x\}$ is convex.} Let
$y,z\in A\setminus\{x\}$, $a,b\ge0$ real with $a+b=1$, and $w=ay+bz$; then $w\in A$. If $w=x$: for $b=0$ we get
$w=y\ne x$, for $a=0$ we get $w=z\ne x$, and for $a,b>0$ extremality gives $y=x$; each case is a contradiction. So
$w\in A\setminus\{x\}$.

Step 2. The disc $B_r$ is convex: for $y,z\in B_r$ and $a,b\ge0$ with $a+b=1$, $\abs{ay+bz}\le a\abs y+b\abs
z\le r$. Let $s\subseteq\iota$ and $x\in\iota$ with $f(x)\in\conv f(s)$, and suppose $x\notin s$. For $j\in s$ we have
$\abs{f(j)}=r$, so $f(j)\in B_r$, and $j\ne x$, so $f(j)\ne f(x)$ by injectivity. Thus $f(s)\subseteq
B_r\setminus\{f(x)\}$. Since $\abs{f(x)}=r$, $f(x)$ is an extreme point of $B_r$ (Lemma~\ref{ec:lem:sphere}), so
$B_r\setminus\{f(x)\}$ is convex by Step~1, and therefore contains $\conv f(s)$. This contradicts $f(x)\in\conv f(s)$.
Hence $x\in s$.

(The Lean proof uses Mathlib's reformulation of convex independence, ``$f(i)\notin\conv f(s\setminus\{i\})$ for all
$i,s$'', with the same two steps.)
\end{proof}

\begin{definition}[the generator; \leanref{RefuteC6.zeta}{ES8SYM/lean/C6check.lean:60}]\label{ec:def:zeta}
For $m\in\N$, $\zeta_m=\exp(2\pi i/m)\in\C$, with $m$ cast to $\C$. (For $m=0$, Lean's division gives $2\pi i/0=0$ and
$\zeta_0=1$; this case is not used.) Rotation by $2\pi/m$ about $0$ is $z\mapsto\zeta_mz$.
\end{definition}

\begin{lemma}[primitive root; \leanref{RefuteC6.zeta_prim}{ES8SYM/lean/C6check.lean:62}]\label{ec:lem:zeta}
For $m\in\N$ with $m\ne0$, $\zeta_m$ is a primitive $m$-th root of unity.
\end{lemma}

\begin{proof}
This is Mathlib's \lean{Complex.isPrimitiveRoot_exp}, a standard fact. Proof: $\exp(w)=1$ if and only if
$w\in2\pi i\Z$ (periodicity of the exponential function), and $\zeta_m^l=\exp(2\pi il/m)$ for $l\in\N$. Hence
$\zeta_m^l=1$ if and only if $l/m\in\Z$, that is $m\mid l$; in particular $\zeta_m^m=1$.
\end{proof}

\begin{fact}[properties of primitive roots; written step]\label{ec:fact:prim}
Let $m\ge1$ and let $\zeta\in\C$ be a primitive $m$-th root of unity. Then
\textup{(i)} $\zeta\ne0$; \textup{(ii)} $\abs\zeta=1$; \textup{(iii)} if $i,j<m$ and $\zeta^i=\zeta^j$, then $i=j$;
\textup{(iv)} $\zeta^j\ne1$ for $0<j<m$.
\end{fact}

\begin{proof}
(i) $\zeta^m=1\ne0$ and $m\ge1$. (ii) $\abs\zeta^m=\abs{\zeta^m}=1$ with $\abs\zeta\ge0$, and $x\mapsto x^m$ is
strictly increasing on $[0,\infty)$, so $\abs\zeta=1$. (iii) Say $i\le j$. Then
$\zeta^i\zeta^{j-i}=\zeta^j=\zeta^i$, and $\zeta^i\ne0$ by (i), so $\zeta^{j-i}=1$, hence $m\mid j-i$; as $0\le
j-i<m$, $j=i$. (iv) $\zeta^j=1$ would give $m\mid j$, impossible for $0<j<m$. (In Mathlib these are
\lean{IsPrimitiveRoot.ne_zero}, \lean{IsPrimitiveRoot.norm'_eq_one}, \lean{IsPrimitiveRoot.pow_inj} and
\lean{IsPrimitiveRoot.pow_ne_one_of_pos_of_lt}.)
\end{proof}

\begin{lemma}[eight points of an orbit; \leanref{RefuteC6.orbit_injective}{ES8SYM/lean/C6check.lean:66}, \leanref{RefuteC6.orbit_convexIndependent}{ES8SYM/lean/C6check.lean:75}]\label{ec:lem:orbit8}
Let $m\in\N$ with $8\le m$ and $p\in\C$ with $p\ne0$, and let $F:\mathrm{Fin}\,8\to\C$, $F(k)=p\,\zeta_m^k$. Then
\textup{(a)} $F$ is injective, and \textup{(b)} $F$ is convex independent over $\R$.
\end{lemma}

\begin{proof}
Since $m\ne0$, $\zeta_m$ is a primitive $m$-th root of unity (Lemma~\ref{ec:lem:zeta}).
(a) If $p\zeta_m^j=p\zeta_m^k$, then $\zeta_m^j=\zeta_m^k$ because $p\ne0$, and $j,k<8\le m$, so $j=k$ by
Fact~\ref{ec:fact:prim}(iii).
(b) $\abs{F(k)}=\abs p\,\abs{\zeta_m}^k=\abs p$ by Fact~\ref{ec:fact:prim}(ii). Apply Lemma~\ref{ec:lem:ci-circle}
with $r=\abs p$ and (a).
\end{proof}

\begin{theorem}[order at least $8$; \leanref{RefuteC6.part_a}{ES8SYM/lean/C6check.lean:85}]\label{ec:thm:big-order}
Let $m\in\N$ with $8\le m$, and let $S\subseteq\C$ be any set \textup{(}not necessarily finite\textup{)} such that
$\exp(2\pi i/m)\,z\in S$ for every $z\in S$. Assume that for every injective $f:\mathrm{Fin}\,8\to\C$ with $f(k)\in S$
for all $k$, $f$ is not convex independent over $\R$. Then $S\subseteq\{0\}$.
\end{theorem}

\begin{proof}
We follow the Lean proof. Let $p\in S$ and suppose $p\ne0$. By induction, $p\zeta_m^k\in S$ for every $k\in\N$:
$p\zeta_m^0=p\in S$, and if $p\zeta_m^k\in S$, then $\exp(2\pi i/m)\,(p\zeta_m^k)=p\zeta_m^{k+1}\in S$ by the hypothesis
(note $\exp(2\pi i/m)=\zeta_m$ by definition). Hence the family $F$ of Lemma~\ref{ec:lem:orbit8} takes values in $S$; it
is injective and convex independent by that lemma, contradicting the assumption. So $p=0$.
\end{proof}

\begin{remark}\label{ec:rem:big-order}
The theorem uses no finiteness and no general position, and only the forward invariance $\zeta_mS\subseteq S$. Its
conclusion says $S=\emptyset$ or $S=\{0\}$; both sets satisfy the hypotheses, since an injective family of $8$ values
needs $8$ distinct points. The planar statement used in the project (\dirn{ES8SYM/proofs/layers.md}, Lemma~2(a),
stated there for sets in general position; the row ``order $\ge8$'', labelled Lean-checked, of the $k=8$ tables in
\dirn{ES8SYM/README.md} and \dirn{WRITEUP/note/symmetric_extremal.tex}) follows from the following corollary, which
does not assume general position. Lemma~2(a) is stated for the rotation by $2\pi/m$, and the corollary applies to it
directly. The tables concern any rotation $\rho$ of order $m$ (\dirn{WRITEUP/note/symmetric_extremal.tex}, line~338),
that is $\rho=\mathrm{rotAbout}(c,2\pi k/m)$ with $\gcd(k,m)=1$; as in Remark~\ref{ec:rem:es7-sharp}, the power
$\rho^j$ with $j\in\N$ and $jk\equiv1\pmod m$ is $\mathrm{rotAbout}(c,2\pi/m)$ and maps every $\rho$-invariant set
into itself, so the corollary applies to $\rho^j$. Lean checks only the complex form, Theorem~\ref{ec:thm:big-order};
the passage to the plane is not formalized in Lean, and a written proof follows.
\end{remark}

\begin{corollary}[planar form of Theorem~\ref{ec:thm:big-order}; written step]\label{ec:cor:planar}
Let $m\in\N$ with $m\ge8$, let $c$ be a point, and let $S$ be a finite set of points that is invariant under
$\mathrm{rotAbout}(c,2\pi/m)$ and satisfies $\neg\,\mathrm{HasConvexNGon}(8,S)$. Then $S\subseteq\{c\}$.
\end{corollary}

\begin{proof}
Let $\varphi:\R^2\to\C$, $\varphi(p)=(p_0-c_0)+i(p_1-c_1)$. It is a bijection and an affine map of real vector spaces,
with inverse $\psi(z)=(c_0+\operatorname{Re}z,\;c_1+\operatorname{Im}z)$. By the complex form of
Definition~\ref{fd:def:rot}, $\varphi(\mathrm{rotAbout}(c,\theta)\,p)=e^{i\theta}\varphi(p)$ for all $\theta\in\R$ and
all points $p$. We also use the standard fact that an affine map $\alpha$ between real vector spaces satisfies
$\alpha(\conv Y)=\conv\alpha(Y)$, because $\conv Y$ is the set of finite convex combinations of points of $Y$ and
$\alpha$ preserves convex combinations. Let $\tilde S=\varphi(S)\subseteq\C$.

\emph{Invariance.} Let $z=\varphi(p)$ with $p\in S$. With $\theta=2\pi/m$ we have $e^{i\theta}=\exp(2\pi i/m)$, so
$\exp(2\pi i/m)\,z=\varphi(\mathrm{rotAbout}(c,2\pi/m)\,p)\in\tilde S$, because $\mathrm{rotAbout}(c,2\pi/m)\,p\in S$.

\emph{No convex independent $8$-family.} Let $f:\mathrm{Fin}\,8\to\C$ be injective with values in $\tilde S$, and
suppose $f$ is convex independent. Put $g=\psi\circ f$; it is injective with values in $S$, so $T=g(\mathrm{Fin}\,8)$
is an $8$-element subset of $S$. We show that $T$ is in convex position, i.e.\ (Definition~\ref{fd:def:cp}) that for
every $U\subseteq T$ and every $x\in T$, $x\in\conv U$ implies $x\in U$. Write $x=g(k)$ and
$s=\{j\in\mathrm{Fin}\,8: g(j)\in U\}$; then $U=g(s)$, and $\varphi\circ g=f$. From $x\in\conv U$ we get
$f(k)=\varphi(x)\in\varphi(\conv U)=\conv\varphi(U)=\conv f(s)$, so $k\in s$ by convex independence of $f$, that is
$x=g(k)\in U$. Hence $\mathrm{HasConvexNGon}(8,S)$, a contradiction.

So $\tilde S$ satisfies the hypotheses of Theorem~\ref{ec:thm:big-order}, and $\tilde S\subseteq\{0\}$. Since
$\varphi(p)=0$ only for $p=c$, $S\subseteq\{c\}$.
\end{proof}

\begin{proposition}[threshold arithmetic; \leanref{RefuteC6.thr_m4}{ES8SYM/lean/C6check.lean:108}, \leanref{RefuteC6.thr_m5}{ES8SYM/lean/C6check.lean:109}, \leanref{RefuteC6.thr_m6}{ES8SYM/lean/C6check.lean:110}, \leanref{RefuteC6.thr_m7}{ES8SYM/lean/C6check.lean:111}, \leanref{RefuteC6.m4_with_centre_counter}{ES8SYM/lean/C6check.lean:115}]\label{ec:prop:thr}
In $\N$:
\begin{enumerate}[label=(\roman*)]
\item \textup{(\lean{thr_m4})} for every $t$: $65\le4t\iff17\le t$;
\item \textup{(\lean{thr_m5})} for all $t,e$ with $e\le1$: $65\le5t+e\iff13\le t$;
\item \textup{(\lean{thr_m6})} for every $t$: $65\le6t\iff11\le t$;
\item \textup{(\lean{thr_m7})} for all $t,e$ with $e\le1$: $65\le7t+e\iff10\le t$;
\item \textup{(\lean{m4_with_centre_counter})} $65\le4\cdot16+1$ and $\neg(17\le16)$.
\end{enumerate}
\textup{(}In Lean the variable $e$ is called \texttt{c}.\textup{)}
\end{proposition}

\begin{proof}
(i) If $t\ge17$ then $4t\ge68$; if $t\le16$ then $4t\le64$. (ii) If $t\ge13$ then $5t+e\ge65$; if $t\le12$ then
$5t+e\le60+1=61$. (iii) If $t\ge11$ then $6t\ge66$; if $t\le10$ then $6t\le60$. (iv) If $t\ge10$ then $7t+e\ge70$; if
$t\le9$ then $7t+e\le63+1=64$. (v) $4\cdot16+1=65$, and $17>16$. (Lean proves each line with \lean{omega}, a decision
procedure for linear arithmetic over $\N$ and $\Z$ whose proof term is checked by the kernel.)
\end{proof}

\begin{remark}[the geometric part of Lemma~2(b), not in this file]\label{ec:rem:thr}
In \dirn{ES8SYM/proofs/layers.md}, Lemma~2(b), the proposition is applied to a finite set $S$ in general position,
invariant under $\mathrm{rotAbout}(c,2\pi/m)$ with $m\in\{4,5,6,7\}$ and without $8$ points in convex position. The
set $S\setminus\{c\}$ is again invariant (the rotation is injective and fixes $c$), and the layer results applied to it
(Theorems~\ref{ll:thm:layers} and~\ref{ll:thm:card-peel}, with $K=8\le2m$) give $\card{S\setminus\{c\}}=mt$, where
$t$ is the number of convex layers of $S\setminus\{c\}$. If moreover $S\not\subseteq\{c\}$ (for instance if
$\card S\ge2$), then $c\notin S$ for even $m$ (Lemma~\ref{lg:lem:centre}). So $\card S=mt+e$ with $e=[c\in S]\le1$,
and $e=0$ for $m=4,6$ when $\card S\ge2$. Then Proposition~\ref{ec:prop:thr} turns $\card S\ge65$ into
$t\ge17,13,11,10$. This geometric input is not part of \dirn{ES8SYM/lean/C6check.lean}, so Lemma~2(b) as a whole is
not Lean-checked there. Item (v) records the arithmetic behind the
remark of \dirn{ES8SYM/proofs/layers.md} (lines 41--42) that general position is needed: with the centre allowed for
$m=4$, $16$ orbits and the centre would already give $65$ points.
\end{remark}

%---------------------------------------------------------------------------------------------------------------------
\subsection{Orientations in inverse-radius polar form}\label{ec:ssec:polar}

The source is \dirn{ES8SYM/lean/PolarOrient.lean} (namespace \lean{C7}; Mathlib only). It is the Lean check of
Lemma~1 of \dirn{ES8SYM/proofs/layers.md} (lines 15--23) and of part of the consequence stated after it (lines 25--29;
see Remark~\ref{ec:rem:polar}(iii)). For $x\in\R$, $\operatorname{sign}x\in\{-1,0,1\}$ is
Mathlib's \lean{SignType.sign}: $1$ if $x>0$, $0$ if $x=0$, $-1$ if $x<0$. We use
$\operatorname{sign}(xy)=\operatorname{sign}(x)\operatorname{sign}(y)$ and $\operatorname{sign}(x)=1$ for $x>0$; hence
$\operatorname{sign}(xP)=\operatorname{sign}(x)$ whenever $P>0$. For $p\in\R\times\R$ we write $p=(p_1,p_2)$ (Lean
\lean{p.1}, \lean{p.2}).

\begin{definition}[cross product, inverse-radius points, layered coefficients; \leanref{C7.cross}{ES8SYM/lean/PolarOrient.lean:45}, \leanref{C7.pt}{ES8SYM/lean/PolarOrient.lean:48}, \leanref{C7.coef}{ES8SYM/lean/PolarOrient.lean:63}, \leanref{C7.phase}{ES8SYM/lean/PolarOrient.lean:94}, \leanref{C7.ngonPt}{ES8SYM/lean/PolarOrient.lean:97}]\label{ec:def:polar}
The following objects are used in this subsection.
\begin{enumerate}[label=(\roman*)]
\item For $p,q,r\in\R\times\R$:
  $\mathrm{cross}(p,q,r)=(q_1-p_1)(r_2-p_2)-(q_2-p_2)(r_1-p_1)$.
\item For $\varphi,u\in\R$: $\mathrm{pt}(\varphi,u)=(\cos\varphi/u,\ \sin\varphi/u)\in\R\times\R$, the point with
  direction $\varphi$ and inverse radius $u$ (Lean's division gives $x/0=0$; only $u>0$ is used).
\item For a type $\iota$, $t\in\N$, $L:\iota\to\mathrm{Fin}\,t$, $\varphi:\iota\to\R$ and $i,j,k\in\iota$, the vector
  $\mathrm{coef}_{L,\varphi}(i,j,k)\in\R^{\mathrm{Fin}\,t}$ is
  \begin{equation*}
  \mathrm{coef}_{L,\varphi}(i,j,k)=\sin(\varphi_k-\varphi_j)\,e_{L(i)}+\sin(\varphi_i-\varphi_k)\,e_{L(j)}
  +\sin(\varphi_j-\varphi_i)\,e_{L(k)},
  \end{equation*}
  where $\varphi_i=\varphi(i)$ and $e_o$ is the $o$-th unit vector (Lean \lean{Pi.single}); when two of
  $L(i),L(j),L(k)$ coincide, the coefficients add.
\item For $m,t\in\N$, $\theta:\mathrm{Fin}\,t\to\R$ and $(o,k)\in\mathrm{Fin}\,t\times\mathrm{Fin}\,m$:
  $\mathrm{phase}_{m,\theta}(o,k)=\theta_o+2\pi k/m$, computed in $\R$ as $(2\pi k)/m$ with $k$ and $m$ cast to $\R$.
\item For moreover $r:\mathrm{Fin}\,t\to\R$: $\mathrm{ngonPt}_{m,\theta,r}(o,k)=\bigl(r_o\cos\mathrm{phase}_{m,\theta}(o,k),\
  r_o\sin\mathrm{phase}_{m,\theta}(o,k)\bigr)$, vertex $k$ of the regular $m$-gon with centre $0$, radius $r_o$ and
  phase $\theta_o$.
\end{enumerate}
\end{definition}

For the rest of this subsection, for real $\varphi_1,\varphi_2,\varphi_3,u_1,u_2,u_3$ let
\begin{equation*}
D=\det\begin{pmatrix}\cos\varphi_1/u_1&\sin\varphi_1/u_1&1\\ \cos\varphi_2/u_2&\sin\varphi_2/u_2&1\\
\cos\varphi_3/u_3&\sin\varphi_3/u_3&1\end{pmatrix},\qquad
\Sigma=u_1\sin(\varphi_3-\varphi_2)+u_2\sin(\varphi_1-\varphi_3)+u_3\sin(\varphi_2-\varphi_1).
\end{equation*}

\begin{lemma}[the value identity; \leanref{C7.det_mul_eq}{ES8SYM/lean/PolarOrient.lean:24}]\label{ec:lem:det-mul}
For all real $\varphi_1,\varphi_2,\varphi_3,u_1,u_2,u_3$ with $u_1\ne0$, $u_2\ne0$, $u_3\ne0$: $D\cdot(u_1u_2u_3)=\Sigma$.
\end{lemma}

\begin{proof}
Write $x_k=\cos\varphi_k/u_k$ and $y_k=\sin\varphi_k/u_k$. Expanding the determinant along its third column,
\begin{equation*}
D=(x_2y_3-x_3y_2)+(x_3y_1-x_1y_3)+(x_1y_2-x_2y_1).
\end{equation*}
For indices $a\ne b$, the subtraction formula $\sin(\beta-\alpha)=\sin\beta\cos\alpha-\cos\beta\sin\alpha$ gives
$x_ay_b-x_by_a=(\cos\varphi_a\sin\varphi_b-\cos\varphi_b\sin\varphi_a)/(u_au_b)=\sin(\varphi_b-\varphi_a)/(u_au_b)$.
Multiplying by $u_1u_2u_3$ turns the three terms into $u_1\sin(\varphi_3-\varphi_2)$, $u_2\sin(\varphi_1-\varphi_3)$
and $u_3\sin(\varphi_2-\varphi_1)$. (The Lean proof expands the determinant by \lean{Matrix.det_fin_three}, applies
the subtraction formula, clears denominators and normalizes; it is the same computation.)
\end{proof}

\begin{lemma}[polar orientation formula; \leanref{C7.orient_polar}{ES8SYM/lean/PolarOrient.lean:37}]\label{ec:lem:polar-u}
For all real $\varphi_1,\varphi_2,\varphi_3,u_1,u_2,u_3$ with $u_1>0$, $u_2>0$, $u_3>0$:
$\operatorname{sign}D=\operatorname{sign}\Sigma$.
\end{lemma}

\begin{proof}
Let $P=u_1u_2u_3>0$. By Lemma~\ref{ec:lem:det-mul}, $\operatorname{sign}\Sigma=\operatorname{sign}(DP)=\operatorname{sign}D$.
\end{proof}

\begin{lemma}[cross-product form; \leanref{C7.orient_polar_cross}{ES8SYM/lean/PolarOrient.lean:50}]\label{ec:lem:polar-cross}
For all real $\varphi_1,\varphi_2,\varphi_3,u_1,u_2,u_3$ with $u_1>0$, $u_2>0$, $u_3>0$:
\begin{equation*}
\operatorname{sign}\mathrm{cross}\bigl(\mathrm{pt}(\varphi_1,u_1),\mathrm{pt}(\varphi_2,u_2),\mathrm{pt}(\varphi_3,u_3)\bigr)
=\operatorname{sign}\Sigma .
\end{equation*}
\end{lemma}

\begin{proof}
With $x_k,y_k$ as before, $\mathrm{pt}(\varphi_k,u_k)=(x_k,y_k)$, and expanding,
\begin{equation*}
(x_2-x_1)(y_3-y_1)-(y_2-y_1)(x_3-x_1)=(x_2y_3-x_3y_2)+(x_3y_1-x_1y_3)+(x_1y_2-x_2y_1),
\end{equation*}
the same three-term expression as $D$. So $\mathrm{cross}(\dots)\cdot u_1u_2u_3=\Sigma$ by the computation in the proof
of Lemma~\ref{ec:lem:det-mul}, and the sign argument of Lemma~\ref{ec:lem:polar-u} applies.
\end{proof}

\begin{proposition}[layered configurations; \leanref{C7.orient_layered}{ES8SYM/lean/PolarOrient.lean:68}]\label{ec:prop:layered}
Let $\iota$ be a type, $t\in\N$, $L:\iota\to\mathrm{Fin}\,t$, $\varphi:\iota\to\R$, and $u:\mathrm{Fin}\,t\to\R$ with
$u(o)>0$ for every $o$. Then for all $i,j,k\in\iota$ \textup{(}not necessarily distinct\textup{)}:
\begin{equation*}
\operatorname{sign}\mathrm{cross}\bigl(\mathrm{pt}(\varphi_i,u(L(i))),\mathrm{pt}(\varphi_j,u(L(j))),\mathrm{pt}(\varphi_k,u(L(k)))\bigr)
=\operatorname{sign}\sum_{o\in\mathrm{Fin}\,t}\mathrm{coef}_{L,\varphi}(i,j,k)_o\,u(o).
\end{equation*}
\end{proposition}

\begin{proof}
For $a\in\mathrm{Fin}\,t$ and $v\in\R$, $\sum_o(v\,e_a)_o\,u(o)=v\,u(a)$, and the sum is linear in the coefficient
vector, so
\begin{equation*}
\sum_o\mathrm{coef}_{L,\varphi}(i,j,k)_o\,u(o)=u(L(i))\sin(\varphi_k-\varphi_j)+u(L(j))\sin(\varphi_i-\varphi_k)
+u(L(k))\sin(\varphi_j-\varphi_i),
\end{equation*}
also when some of $L(i),L(j),L(k)$ coincide. This is $\Sigma$ for $(\varphi_1,\varphi_2,\varphi_3)=(\varphi_i,\varphi_j,\varphi_k)$
and $(u_1,u_2,u_3)=(u(L(i)),u(L(j)),u(L(k)))$, all positive; apply Lemma~\ref{ec:lem:polar-cross}.
\end{proof}

\begin{corollary}[scaling; \leanref{C7.orient_layered_scale}{ES8SYM/lean/PolarOrient.lean:80}]\label{ec:cor:scale}
In the setting of Proposition~\ref{ec:prop:layered}, let $\lambda\in\R$ with $\lambda>0$. Then for all $i,j,k\in\iota$,
\begin{multline*}
\operatorname{sign}\mathrm{cross}\bigl(\mathrm{pt}(\varphi_i,\lambda u(L(i))),\mathrm{pt}(\varphi_j,\lambda u(L(j))),\mathrm{pt}(\varphi_k,\lambda u(L(k)))\bigr)\\
=\operatorname{sign}\mathrm{cross}\bigl(\mathrm{pt}(\varphi_i,u(L(i))),\mathrm{pt}(\varphi_j,u(L(j))),\mathrm{pt}(\varphi_k,u(L(k)))\bigr).
\end{multline*}
\end{corollary}

\begin{proof}
The function $o\mapsto\lambda u(o)$ is positive, so Proposition~\ref{ec:prop:layered} applies to it and to $u$. With
$c_o=\mathrm{coef}_{L,\varphi}(i,j,k)_o$, the left side is
$\operatorname{sign}\sum_oc_o\lambda u(o)=\operatorname{sign}\bigl(\lambda\sum_oc_ou(o)\bigr)=\operatorname{sign}\sum_oc_ou(o)$,
which is the right side. (Directly: $\mathrm{pt}(\varphi,\lambda u)=\lambda^{-1}\mathrm{pt}(\varphi,u)$ and
$\mathrm{cross}$ is homogeneous of degree $2$.)
\end{proof}

\begin{lemma}[polygon vertices in inverse-radius form; \leanref{C7.ngonPt_eq_pt}{ES8SYM/lean/PolarOrient.lean:100}]\label{ec:lem:ngon-pt}
Let $m,t\in\N$, $\theta,r:\mathrm{Fin}\,t\to\R$ with $r_o>0$ for every $o$, and $p\in\mathrm{Fin}\,t\times\mathrm{Fin}\,m$.
Then $\mathrm{ngonPt}_{m,\theta,r}(p)=\mathrm{pt}\bigl(\mathrm{phase}_{m,\theta}(p),\,1/r_{p_1}\bigr)$, where $p_1$ is
the first component of $p$.
\end{lemma}

\begin{proof}
For $r_o\ne0$ and every real $\phi$, $\cos\phi/(1/r_o)=r_o\cos\phi$ and $\sin\phi/(1/r_o)=r_o\sin\phi$.
\end{proof}

\begin{theorem}[nested regular $m$-gons; \leanref{C7.nested_regular_orient}{ES8SYM/lean/PolarOrient.lean:108}]\label{ec:thm:nested}
Let $m,t\in\N$ and $\theta,r:\mathrm{Fin}\,t\to\R$ with $r_o>0$ for every $o$. Then for all
$a,b,c\in\mathrm{Fin}\,t\times\mathrm{Fin}\,m$:
\begin{equation*}
\operatorname{sign}\mathrm{cross}\bigl(\mathrm{ngonPt}_{m,\theta,r}(a),\mathrm{ngonPt}_{m,\theta,r}(b),\mathrm{ngonPt}_{m,\theta,r}(c)\bigr)
=\operatorname{sign}\sum_{o\in\mathrm{Fin}\,t}\mathrm{coef}_{\pi_1,\mathrm{phase}_{m,\theta}}(a,b,c)_o\cdot\frac1{r_o},
\end{equation*}
where $\pi_1:\mathrm{Fin}\,t\times\mathrm{Fin}\,m\to\mathrm{Fin}\,t$ is the first projection \textup{(}Lean
\lean{Prod.fst}\textup{)}. The coefficient vector depends only on $m$ and $\theta$, not on $r$.
\end{theorem}

\begin{proof}
By Lemma~\ref{ec:lem:ngon-pt}, $\mathrm{ngonPt}_{m,\theta,r}(a)=\mathrm{pt}(\mathrm{phase}_{m,\theta}(a),u(\pi_1(a)))$
with $u(o)=1/r_o>0$, and likewise for $b$ and $c$. Apply Proposition~\ref{ec:prop:layered} with
$\iota=\mathrm{Fin}\,t\times\mathrm{Fin}\,m$, $L=\pi_1$, $\varphi=\mathrm{phase}_{m,\theta}$ and this $u$.
\end{proof}

\begin{lemma}[triples with the centre; \leanref{C7.orient_centre}{ES8SYM/lean/PolarOrient.lean:116}]\label{ec:lem:centre-u}
For all real $\varphi_1,\varphi_2,u_1,u_2$ with $u_1>0$ and $u_2>0$:
$\operatorname{sign}\mathrm{cross}\bigl(\mathrm{pt}(\varphi_1,u_1),\mathrm{pt}(\varphi_2,u_2),(0,0)\bigr)=\operatorname{sign}\sin(\varphi_2-\varphi_1)$.
\end{lemma}

\begin{proof}
With $x_k,y_k$ as before, $\mathrm{cross}((x_1,y_1),(x_2,y_2),(0,0))=(x_2-x_1)(-y_1)-(y_2-y_1)(-x_1)=x_1y_2-x_2y_1
=\sin(\varphi_2-\varphi_1)/(u_1u_2)$. Multiply by $u_1u_2>0$ and argue as in Lemma~\ref{ec:lem:polar-u}.
\end{proof}

\begin{remark}[scope of the polar lemmas]\label{ec:rem:polar}
(i) If $\hat p$ denotes the point with coordinates $\hat p_0=p_1$, $\hat p_1=p_2$, then
$\mathrm{cross}(p,q,r)=\orient(\hat p,\hat q,\hat r)$ by the formulas (Definition~\ref{fd:def:orient}), so a positive
value means counterclockwise; Lean has no link between \lean{C7.cross} and \lean{ErdosSzekeres.orient}. By the
expansions in the proofs of Lemmas~\ref{ec:lem:det-mul} and~\ref{ec:lem:polar-cross}, $\mathrm{cross}(p,q,r)$ is
the determinant of the matrix with rows $(p_1,p_2,1)$, $(q_1,q_2,1)$, $(r_1,r_2,1)$.
(ii) Lemma~\ref{ec:lem:polar-u} is not the theorem \lean{orient_polar} of \dirn{LAYERLEAN/LayerLemma.lean:935}
(Lemma~\ref{lg:lem:polar}), which is a radius form about a centre $c$:
$\orient(\mathrm{polar}(c,r_a,\alpha),\mathrm{polar}(c,r_b,\beta),\mathrm{polar}(c,r_d,\delta))=r_ar_b\sin(\beta-\alpha)
+r_br_d\sin(\delta-\beta)+r_dr_a\sin(\alpha-\delta)$, where $\mathrm{polar}(c,r,\alpha)=(c_0+r\cos\alpha,\,c_1+r\sin\alpha)$.
For $u\ne0$ the point $\mathrm{pt}(\varphi,u)$ corresponds as in (i) to $\mathrm{polar}(0,1/u,\varphi)$, and for
$r_a=1/u_1$, $r_b=1/u_2$, $r_d=1/u_3$ this right side times $u_1u_2u_3$ is $\Sigma$ with
$(\varphi_1,\varphi_2,\varphi_3)=(\alpha,\beta,\delta)$, so the two agree in substance,
but they are distinct declarations; the name \lean{orient_polar} in \dirn{WRITEUP/P1/p1.tex} and
\dirn{WRITEUP/note/k8_soundness.tex} refers to the radius form.
(iii) Theorem~\ref{ec:thm:nested} assumes only $r_o>0$: no nesting, no order of the radii, no distinctness of the
points and no general position. It says that for fixed phases every orientation sign is the sign of a linear form in
$(1/r_o)_o$ whose coefficients depend only on $\theta$ and $m$; Corollary~\ref{ec:cor:scale} adds that the signs do
not change when all inverse radii are multiplied by the same $\lambda>0$. Lemma~\ref{ec:lem:centre-u} treats triples
that contain the centre $(0,0)$, which have no finite inverse radius. The further statement of
\dirn{ES8SYM/proofs/layers.md} (lines 26--27) that for fixed phases the order type is constant on the cells of a
central hyperplane arrangement is not formalized as such.
\end{remark}

%---------------------------------------------------------------------------------------------------------------------
\subsection{Size compatibility in a group with zero}\label{ec:ssec:compat}

The source is \dirn{PRIMEFOLD/compat/lean/PCompatCore.lean} (namespace \lean{PCompat}). Here the plane is replaced
by an arbitrary commutative group with zero $K$ (Mathlib \lean{CommGroupWithZero}: a commutative monoid with zero in
which $0\ne1$, with an inverse map such that $aa^{-1}=1$ for every $a\ne0$ and $0^{-1}=0$; the structure also carries a
division and integer powers, which are not used here), with decidable equality (Lean needs it to form finite sets by
image, difference, erasure and insertion, and for the test $0\in S$ in Theorem~\ref{ec:thm:compat}(b); mathematically
it is no restriction). In the application $K=\C$, the centre is $0$ and the rotation is
$z\mapsto\zeta z$; a rotation about another centre reduces to this by a translation. We use only these consequences of
the axioms: $0\cdot a=0$; if $a\ne0$ and $ab=ac$, then $b=c$ (multiply by $a^{-1}$); if $a\ne0$ and $b\ne0$, then
$ab\ne0$ (otherwise $b=a^{-1}(ab)=0$); hence $a^j\ne0$ for $a\ne0$. For a finite $T\subseteq K$, ``$T$ is
$\zeta$-invariant'' means $\zeta z\in T$ for every $z\in T$. We write $[P]$ for $1$ if $P$ holds and $0$ otherwise,
and $N\bmod m$ for the remainder (Lean \texttt{N \% m}). The hypotheses on $m\in\N$, $\zeta\in K$ and $r:\N\to K$ are:
\begin{itemize}
\item[(O)] $0<m$, $\zeta^m=1$, and $\zeta^j\ne1$ for every $j$ with $0<j<m$ \ ($\zeta$ has exact order $m$);
\item[(R0)] $r_i\ne0$ for every $i\in\N$;
\item[(RD)] for all $i,i',j,j'\in\N$ with $j<m$ and $j'<m$: if $\zeta^jr_i=\zeta^{j'}r_{i'}$, then $i=i'$ \ (the
  orbits of the $r_i$ are pairwise disjoint).
\end{itemize}

\begin{lemma}[powers; \leanref{PCompat.ne_zero_of_order}{PRIMEFOLD/compat/lean/PCompatCore.lean:36}, \leanref{PCompat.pow_inj_of_order}{PRIMEFOLD/compat/lean/PCompatCore.lean:43}, \leanref{PCompat.pow_mul_mem}{PRIMEFOLD/compat/lean/PCompatCore.lean:58}]\label{ec:lem:pc-basic}
Let $K$ be a commutative group with zero, $m\in\N$ and $\zeta\in K$.
\begin{enumerate}[label=(\alph*)]
\item If $0<m$ and $\zeta^m=1$, then $\zeta\ne0$.
\item If \textup{(O)} holds, $i<m$, $j<m$ and $\zeta^i=\zeta^j$, then $i=j$.
\item If $T\subseteq K$ is finite and $\zeta$-invariant and $z\in T$, then $\zeta^jz\in T$ for every $j\in\N$.
\end{enumerate}
\end{lemma}

\begin{proof}
(a) If $\zeta=0$, then $\zeta^m=0$ since $m>0$, so $0=1$, which is false in $K$.
(b) By (a), $\zeta\ne0$. Suppose $i<j$. Then $\zeta^i\zeta^{j-i}=\zeta^j=\zeta^i=\zeta^i\cdot1$, and $\zeta^i\ne0$, so
$\zeta^{j-i}=1$ with $0<j-i<m$, contradicting (O). The case $j<i$ is symmetric. Hence $i=j$.
(c) Induction on $j$: $\zeta^0z=z\in T$, and $\zeta^{j+1}z=\zeta(\zeta^jz)\in T$ by invariance.
\end{proof}

\begin{lemma}[the converse constructions; \leanref{PCompat.exists_card0}{PRIMEFOLD/compat/lean/PCompatCore.lean:129}, \leanref{PCompat.exists_card1}{PRIMEFOLD/compat/lean/PCompatCore.lean:160}]\label{ec:lem:pc-exist}
Let $K$ be a commutative group with zero, $m\in\N$, $\zeta\in K$ and $r:\N\to K$, and assume \textup{(O)},
\textup{(R0)} and \textup{(RD)}. Then for every $t\in\N$:
\begin{enumerate}[label=(\alph*)]
\item there is a finite $\zeta$-invariant $S\subseteq K$ with $0\notin S$ and $\card S=mt$;
\item there is a finite $\zeta$-invariant $S\subseteq K$ with $\card S=mt+1$.
\end{enumerate}
\end{lemma}

\begin{proof}
(a) Let $P=\{0,\dots,m-1\}\times\{0,\dots,t-1\}$ and $S=\{\zeta^jr_i:(j,i)\in P\}$. \emph{Invariance:} for $(j,i)\in
P$, if $j+1<m$ then $\zeta\cdot\zeta^jr_i=\zeta^{j+1}r_i$ with $(j+1,i)\in P$; if $j+1=m$ then
$\zeta\cdot\zeta^jr_i=\zeta^mr_i=r_i=\zeta^0r_i$ with $(0,i)\in P$ (as $0<m$). \emph{$0\notin S$:} $\zeta\ne0$ by
Lemma~\ref{ec:lem:pc-basic}(a), so $\zeta^j\ne0$, and $r_i\ne0$ by (R0), so $\zeta^jr_i\ne0$. \emph{Cardinality:} the
map $(j,i)\mapsto\zeta^jr_i$ is injective on $P$: if $\zeta^jr_i=\zeta^{j'}r_{i'}$ with $(j,i),(j',i')\in P$, then
$i=i'$ by (RD), then $\zeta^j=\zeta^{j'}$ by cancelling $r_i\ne0$, then $j=j'$ by Lemma~\ref{ec:lem:pc-basic}(b). Hence
$\card S=\card P=mt$.

(b) Take $S$ from (a) and $S'=S\cup\{0\}$. It is $\zeta$-invariant ($\zeta\cdot0=0$), and $\card{S'}=mt+1$ because
$0\notin S$.
\end{proof}

\begin{definition}[size condition; \leanref{PCompat.Compat}{PRIMEFOLD/compat/lean/PCompatCore.lean:205}]\label{ec:def:compat}
For $m,N\in\N$, $\mathrm{Compat}(m,N)$ means: $N\bmod m=0$ or $N\bmod m=1\bmod m$. (For every $m\ge1$ this is the
same condition as $N\bmod m\in\{0,1\}$: for $m\ge2$, $1\bmod m=1$, and for $m=1$ both conditions hold, since
$N\bmod1=0$. The form $1\bmod m$ matches the right side of Theorem~\ref{ec:thm:compat}(b). For $m=0$, Lean's
convention $N\bmod0=N$ applies, but $m=0$ is not used.)
\end{definition}

\begin{proposition}[parity; \leanref{PCompat.compat_even}{PRIMEFOLD/compat/lean/PCompatCore.lean:208}, \leanref{PCompat.compat_odd}{PRIMEFOLD/compat/lean/PCompatCore.lean:235}]\label{ec:prop:parity}
Both parts are statements about natural numbers.
\begin{enumerate}[label=(\alph*)]
\item For every even $m\in\N$ with $2\le m$: $\mathrm{Compat}(m,2^{m-1})$ holds if and only if $m=2^a$ for some
  $a\ge1$.
\item For every odd $m\in\N$ with $3\le m$: $\mathrm{Compat}(m,2^{m-1})$ holds if and only if $2^{m-1}\bmod m=1$.
\end{enumerate}
\end{proposition}

\begin{proof}
We use the standard fact (unique factorization; Mathlib \lean{Nat.dvd_prime_pow}) that the divisors of $2^n$ are
exactly $2^0,2^1,\dots,2^n$. In both parts $m\ge2$, so $1\bmod m=1$ and $\mathrm{Compat}(m,2^{m-1})$ says
$2^{m-1}\bmod m\in\{0,1\}$.

(a) Suppose $2^{m-1}\bmod m=0$. Then $m\mid2^{m-1}$, so $m=2^a$ with $a\le m-1$, and $a\ne0$ because $m\ge2$. Suppose
$2^{m-1}\bmod m=1$. Write $2^{m-1}=qm+1$; since $2\mid m$, $2^{m-1}$ is odd. But $m-1\ge1$, so $2^{m-1}$ is even, a
contradiction. Conversely, let $m=2^a$ with $a\ge1$. Since $a<2^a=m$, we have $a\le m-1$, so $m=2^a\mid2^{m-1}$ and
$2^{m-1}\bmod m=0$.

(b) Suppose $2^{m-1}\bmod m=0$. Then $m=2^a$ for some $a$; $a=0$ gives $m=1<3$, and $a\ge1$ makes $m$ even; both are
impossible. So $\mathrm{Compat}(m,2^{m-1})$ holds exactly when $2^{m-1}\bmod m=1$.
\end{proof}

\begin{theorem}[compatibility; \leanref{PCompat.dvd_card}{PRIMEFOLD/compat/lean/PCompatCore.lean:67}, \leanref{PCompat.card_mod}{PRIMEFOLD/compat/lean/PCompatCore.lean:112}, \leanref{PCompat.exists_iff}{PRIMEFOLD/compat/lean/PCompatCore.lean:174}, \leanref{PCompat.classification}{PRIMEFOLD/compat/lean/PCompatCore.lean:253}]\label{ec:thm:compat}
Let $m\in\N$, and in \textup{(a)}--\textup{(c)} let $K$ be a commutative group with zero and $\zeta\in K$.
\begin{enumerate}[label=(\alph*)]
\item \textup{(\lean{dvd_card})} If \textup{(O)} holds, then every finite $\zeta$-invariant $T\subseteq K$ with
  $0\notin T$ satisfies $m\mid\card T$.
\item \textup{(\lean{card_mod})} If \textup{(O)} holds, then every finite $\zeta$-invariant $S\subseteq K$ satisfies
  $\card S\bmod m=[0\in S]\bmod m$.
\item \textup{(\lean{exists_iff})} Let $r:\N\to K$. If \textup{(O)}, \textup{(R0)} and \textup{(RD)} hold, then for
  every $N\in\N$: there is a finite $\zeta$-invariant $S\subseteq K$ with $\card S=N$ if and only if $N\bmod m=0$ or
  $N\bmod m=1\bmod m$ \textup{(}that is, $\mathrm{Compat}(m,N)$\textup{)}.
\item \textup{(\lean{classification}; this part does not involve $K$ and $\zeta$)} If $1\le m$, then
  $\mathrm{Compat}(m,2^{m-1})$ holds if and only if $m=1$, or $m=2^a$ for some $a\ge1$, or \textup{(}$m$ is odd,
  $3\le m$ and $2^{m-1}\bmod m=1$\textup{)}.
\end{enumerate}
\end{theorem}

\begin{proof}
We follow the Lean proofs.

(a) By Lemma~\ref{ec:lem:pc-basic}(a), $\zeta\ne0$. We use strong induction on $\card T$. If $T=\emptyset$, then
$m\mid0$. Otherwise pick $z\in T$; $z\ne0$ because $0\notin T$. Let $O=\{\zeta^jz:0\le j<m\}$. By
Lemma~\ref{ec:lem:pc-basic}(c), $O\subseteq T$. The map $j\mapsto\zeta^jz$ is injective on $\{0,\dots,m-1\}$: if
$\zeta^iz=\zeta^jz$, then $\zeta^i=\zeta^j$ (cancel $z\ne0$) and $i=j$ by Lemma~\ref{ec:lem:pc-basic}(b). So
$\card O=m$, and $\card T=\card{T\setminus O}+m$. The set $T\setminus O$ omits $0$ and is $\zeta$-invariant: let
$w\in T\setminus O$; then $\zeta w\in T$. Suppose $\zeta w\in O$, say $\zeta w=\zeta^jz$ with $j<m$. If $j=0$, then
$\zeta^{m-1}z=\zeta^{m-1}\zeta w=\zeta^mw=w$, so $w\in O$ (as $m-1<m$). If $j\ge1$, then
$\zeta w=\zeta\cdot\zeta^{j-1}z$, so $w=\zeta^{j-1}z\in O$ (cancel $\zeta\ne0$). Both contradict $w\notin O$, so
$\zeta w\in T\setminus O$. Since $z\in O$, $\card{T\setminus O}<\card T$, and the induction hypothesis gives
$m\mid\card{T\setminus O}$. Hence $m\mid\card T$.

(b) Let $T_0=S\setminus\{0\}$. Then $0\notin T_0$, and $T_0$ is $\zeta$-invariant: for $z\in T_0$, $\zeta z\in S$ and
$\zeta z\ne0$ since $\zeta\ne0$ and $z\ne0$. By (a), $\card{T_0}=mt$ for some $t$. If $0\in S$, then $\card S=mt+1$
and $\card S\bmod m=1\bmod m=[0\in S]\bmod m$. If $0\notin S$, then $S=T_0$ and $\card S\bmod m=0=[0\in S]\bmod m$.

(c) ($\Rightarrow$) Let $S$ be finite, $\zeta$-invariant, with $\card S=N$. By (b), $N\bmod m=1\bmod m$ if $0\in S$,
and $N\bmod m=0\bmod m=0$ if $0\notin S$. ($\Leftarrow$) Write $N=m\lfloor N/m\rfloor+(N\bmod m)$. If $N\bmod m=0$,
Lemma~\ref{ec:lem:pc-exist}(a) with $t=\lfloor N/m\rfloor$ gives a $\zeta$-invariant $S$ with $\card S=N$. Otherwise
$N\bmod m=1\bmod m$. Then $m\ne1$ (for $m=1$ every remainder is $0$), so $m>1$, $1\bmod m=1$, $N=m\lfloor N/m\rfloor+1$,
and Lemma~\ref{ec:lem:pc-exist}(b) with $t=\lfloor N/m\rfloor$ gives the set.

(d) If $m$ is even, then $m\ge2$, and by Proposition~\ref{ec:prop:parity}(a) $\mathrm{Compat}(m,2^{m-1})$ holds iff
$m=2^a$ with $a\ge1$; the other two alternatives of the right side are false for even $m$. If $m$ is odd, then $m=1$ or
$m\ge3$. For $m=1$ both sides hold: $2^0\bmod1=0$, and $m=1$. For odd $m\ge3$, by Proposition~\ref{ec:prop:parity}(b)
the left side holds iff $2^{m-1}\bmod m=1$; on the right side $m=1$ is false and $m=2^a$ with $a\ge1$ is impossible
($m$ is odd), so the right side holds iff $2^{m-1}\bmod m=1$.
\end{proof}

\begin{proposition}[$m=341$; computer-assisted in Lean; written proof by hand; \leanref{PCompat.ex341}{PRIMEFOLD/compat/lean/PCompatCore.lean:277}]\label{ec:prop:ex341}
$\mathrm{Compat}(341,2^{341-1})$; $341=11\cdot31$; $341$ is not prime; $2^{10}\bmod11=1$; $2^{10}\bmod31=1$;
$2^{10}\bmod341=1$.
\end{proposition}

\begin{proof}
This is a finite computation with closed numerals. Lean checks the five numerical claims by kernel evaluation
(\lean{decide}, with \lean{exponentiation.threshold} raised to $400$; no \lean{native_decide}) and non-primality by
Mathlib's \lean{Nat.not_prime_mul} with $11\ne1$ and $31\ne1$; the log lists only \lean{propext}. The computation is
small enough to do by hand, as follows. $11\cdot31=341$, and $11$ is a divisor with
$1<11<341$, so $341$ is not prime. Further $2^{10}=1024=93\cdot11+1=33\cdot31+1=3\cdot341+1$, which gives the three
remainders. Finally $2^{340}=(2^{10})^{34}\equiv1^{34}=1\pmod{341}$, so $2^{340}\bmod341=1=1\bmod341$ and
$\mathrm{Compat}(341,2^{340})$ holds by its second alternative.
\end{proof}

\begin{remark}[what is and is not formalized]\label{ec:rem:compat}
(i) For $K=\C$ and $\zeta=\zeta_m=e^{2\pi i/m}$ with $m\ge1$, (O) holds by Lemma~\ref{ec:lem:zeta} and
Fact~\ref{ec:fact:prim}(iv). (ii) In Theorem~\ref{ec:thm:compat}(c) the hypotheses (R0) and (RD) on $r$ are not
discharged in Lean for $\C$. They hold for $K=\C$, $\zeta=\zeta_m$ and $r_i=i+1$ (a written step): (R0) is clear, and
$\abs{\zeta_m^jr_i}=\abs{\zeta_m}^j(i+1)=i+1$ by Fact~\ref{ec:fact:prim}(ii), so $\zeta_m^jr_i=\zeta_m^{j'}r_{i'}$ forces
$i+1=i'+1$. The sets of Lemma~\ref{ec:lem:pc-exist} need not be in general position (for $m=2$ and $r_i=i+1$ they lie
on the real axis); the existence of size-compatible sets in general position
(\dirn{WRITEUP/note/symmetric_extremal.tex}, the proposition labelled \texttt{prop:compat}) is not formalized.
(iii) Fermat's little theorem is not formalized: the third alternative of Theorem~\ref{ec:thm:compat}(d) is the
arithmetic condition $2^{m-1}\bmod m=1$; by Fermat's little theorem it holds for every odd prime $m$, and for odd
composite $m$ it says that $m$ is a base-$2$ Fermat pseudoprime, as for $m=341$.
\end{remark}

%---------------------------------------------------------------------------------------------------------------------
\subsection{Symmetric sets of \texorpdfstring{$2^{m-1}$}{2\^{}(m-1)} points}\label{ec:ssec:l0}

The source is \dirn{PRIMEFOLD/compat/lean/L0Check.lean} (namespace \lean{L0Check}), an independent check of Lemma~0
of \dirn{PRIMEFOLD/cluster/PROOFS.md}. The plane is $\C$, the centre is $0$, and the rotation by $2\pi/m$ is
multiplication by a primitive $m$-th root of unity. Mathlib's base-$2$ Fermat pseudoprimes are defined by
\begin{equation*}
\mathrm{FermatPsp}(m,2)\iff m\mid2^{m-1}-1,\ m\text{ is not prime, and }1<m
\end{equation*}
(\lean{Nat.FermatPsp}, with natural subtraction; $2^{m-1}\ge1$, so $2^{m-1}-1$ is the usual difference).

\begin{lemma}[arithmetic core; \leanref{L0Check.arith}{PRIMEFOLD/compat/lean/L0Check.lean:21}]\label{ec:lem:arith}
Let $m,t,e\in\N$ with $2\le m$, $e\le1$ and $mt+e=2^{m-1}$. Then either \textup{(}$m$ is odd, $e=1$, and $m$ is prime
or $\mathrm{FermatPsp}(m,2)$\textup{)}, or \textup{(}$e=0$ and $m=2^j$ for some $j\in\N$\textup{)}.
\end{lemma}

\begin{proof}
\emph{$m$ even.} Then $mt$ is even, and $2^{m-1}$ is even because $m-1\ge1$. Hence $e=2^{m-1}-mt$ is even, and $e\le1$
gives $e=0$. Then $mt=2^{m-1}$, so $m\mid2^{m-1}$ and $m=2^j$ for some $j$ (the divisors of $2^n$ are powers of $2$).
The second alternative holds.

\emph{$m$ odd.} If $e\ne1$, then $e=0$, $m\mid2^{m-1}$ and $m=2^j$; $j=0$ gives $m=1<2$, and $j\ge1$ makes $m$ even;
both are impossible. So $e=1$ and $mt+1=2^{m-1}$, hence $m\mid mt=2^{m-1}-1$. If $m$ is prime, the first
alternative holds; otherwise $m\mid2^{m-1}-1$, $m$ is not prime and $1<m$, which is $\mathrm{FermatPsp}(m,2)$.
\end{proof}

\begin{lemma}[orbit counting in $\C$; \leanref{L0Check.orbit_dvd}{PRIMEFOLD/compat/lean/L0Check.lean:60}]\label{ec:lem:orbit-dvd}
Let $m\in\N$ with $0<m$, and let $\zeta\in\C$ be a primitive $m$-th root of unity. Then every finite $T\subseteq\C$
with $0\notin T$ and $\zeta z\in T$ for all $z\in T$ satisfies $m\mid\card T$.
\end{lemma}

\begin{proof}
$\C$ is a commutative group with zero, $\zeta^m=1$, and $\zeta^j\ne1$ for $0<j<m$ by Fact~\ref{ec:fact:prim}(iv); so
(O) holds and Theorem~\ref{ec:thm:compat}(a) applies. (The Lean file proves this statement again, independently, by
the same strong induction, using Mathlib's \lean{IsPrimitiveRoot.pow_inj} for the size of an orbit.)
\end{proof}

\begin{theorem}[sets of $2^{m-1}$ points; \leanref{L0Check.L0}{PRIMEFOLD/compat/lean/L0Check.lean:122}, \leanref{L0Check.L0_rot}{PRIMEFOLD/compat/lean/L0Check.lean:144}]\label{ec:thm:l0}
Let $m\in\N$ with $2\le m$.
\begin{enumerate}[label=(\alph*)]
\item \textup{(\lean{L0})} Let $\zeta\in\C$ be a primitive $m$-th root of unity, and let $S\subseteq\C$ be finite with
  $\zeta z\in S$ for all $z\in S$ and $\card S=2^{m-1}$. Then either \textup{(}$m$ is odd, $0\in S$, and $m$ is prime
  or $\mathrm{FermatPsp}(m,2)$\textup{)}, or \textup{(}$0\notin S$ and $m=2^j$ for some $j\in\N$\textup{)}.
\item \textup{(\lean{L0_rot})} Let $S\subseteq\C$ be finite with $\exp(2\pi i/m)\,z\in S$ for all $z\in S$ and
  $\card S=2^{m-1}$. Then the same conclusion holds.
\end{enumerate}
\end{theorem}

\begin{proof}
(a) By Fact~\ref{ec:fact:prim}(i), $\zeta\ne0$. The set $T=S\setminus\{0\}$ omits $0$ and is invariant: for $z\in T$,
$\zeta z\in S$ and $\zeta z\ne0$. By Lemma~\ref{ec:lem:orbit-dvd}, $\card T=mt$ for some $t\in\N$. If $0\in S$, then
$mt+1=\card S=2^{m-1}$, and Lemma~\ref{ec:lem:arith} with $e=1$ gives the first alternative of the lemma (the second
needs $e=0$); together with $0\in S$ this is the first alternative of the theorem. If $0\notin S$, then $T=S$ and
$mt+0=2^{m-1}$; Lemma~\ref{ec:lem:arith} with $e=0$ excludes its first alternative (which needs $e=1$), so $m=2^j$ for
some $j$, and with $0\notin S$ this is the second alternative of the theorem.

(b) $\exp(2\pi i/m)=\zeta_m$ is a primitive $m$-th root of unity by Lemma~\ref{ec:lem:zeta} ($m\ge2>0$); apply (a).
\end{proof}

\begin{remark}[on the hypotheses]\label{ec:rem:l0}
No general-position or convexity hypothesis is assumed. For even $m$ the centre is excluded by parity, as in the proof
of Lemma~\ref{ec:lem:arith}: $\card S=2^{m-1}$ is even and $\card{S\setminus\{0\}}$ is a multiple of $m$. (Lemma~0 of
\dirn{PRIMEFOLD/cluster/PROOFS.md} excludes the centre by general position instead; the Lean statement does not
need it.) For odd $m$ the theorem only says that $m$ is prime or a base-$2$ Fermat pseudoprime; that every odd prime
actually satisfies $2^{m-1}\equiv1\pmod m$ is Fermat's little theorem, which is not used and not formalized here. For
a rotation about a centre $c\ne0$, translate by $-c$ first. For $k=8$, i.e.\ $m=7$ and $2^{m-1}=64=7\cdot9+1$: a
finite $S\subseteq\C$ with $e^{2\pi i/7}S\subseteq S$ and $\card S=64$ contains $0$, by part (b), since $7$ is odd and
not a power of $2$.
\end{remark}

%---------------------------------------------------------------------------------------------------------------------
\paragraph{Coverage.} Covered, by label:
\begin{itemize}
\item \dirn{LEANEXT/es7corr/ES7Corollaries.lean}, Section~11: \lean{Fivefold35} and \lean{Fourfold32}
  (Definition~\ref{ec:def:cases}); \lean{fivefold_heptagon_of_erase} (Theorem~\ref{ec:thm:fivefold-erase});
  \lean{fivefold_heptagon} (Theorem~\ref{ec:thm:fivefold}); \lean{fivefold_heptagon_of_centre_not_mem}
  (Theorem~\ref{ec:thm:fivefold-nc}); \lean{fourfold_heptagon_of_erase} (Theorem~\ref{ec:thm:fourfold-erase});
  \lean{fourfold_heptagon} (Theorem~\ref{ec:thm:fourfold}). The three fivefold theorems are conditional on
  \lean{Fivefold35}, the two fourfold theorems on \lean{Fourfold32}.
\item \dirn{ES8SYM/lean/C6check.lean}: \lean{sphere_subset_extremePoints} (Lemma~\ref{ec:lem:sphere});
  \lean{convexIndependent_of_norm_eq} (Lemma~\ref{ec:lem:ci-circle}); \lean{zeta} (Definition~\ref{ec:def:zeta});
  \lean{zeta_prim} (Lemma~\ref{ec:lem:zeta}); \lean{orbit_injective} and \lean{orbit_convexIndependent}
  (Lemma~\ref{ec:lem:orbit8}); \lean{part_a} (Theorem~\ref{ec:thm:big-order}); \lean{thr_m4}--\lean{thr_m7} and
  \lean{m4_with_centre_counter} (Proposition~\ref{ec:prop:thr}).
\item \dirn{ES8SYM/lean/PolarOrient.lean}: \lean{cross}, \lean{pt}, \lean{coef}, \lean{phase}, \lean{ngonPt}
  (Definition~\ref{ec:def:polar}); \lean{det_mul_eq} (Lemma~\ref{ec:lem:det-mul}); \lean{orient_polar}
  (Lemma~\ref{ec:lem:polar-u}); \lean{orient_polar_cross} (Lemma~\ref{ec:lem:polar-cross}); \lean{orient_layered}
  (Proposition~\ref{ec:prop:layered}); \lean{orient_layered_scale} (Corollary~\ref{ec:cor:scale});
  \lean{ngonPt_eq_pt} (Lemma~\ref{ec:lem:ngon-pt}); \lean{nested_regular_orient} (Theorem~\ref{ec:thm:nested});
  \lean{orient_centre} (Lemma~\ref{ec:lem:centre-u}).
\item \dirn{PRIMEFOLD/compat/lean/PCompatCore.lean}: \lean{ne_zero_of_order}, \lean{pow_inj_of_order},
  \lean{pow_mul_mem} (Lemma~\ref{ec:lem:pc-basic}); \lean{exists_card0}, \lean{exists_card1}
  (Lemma~\ref{ec:lem:pc-exist}); \lean{Compat} (Definition~\ref{ec:def:compat}); \lean{compat_even},
  \lean{compat_odd} (Proposition~\ref{ec:prop:parity}); \lean{dvd_card}, \lean{card_mod}, \lean{exists_iff},
  \lean{classification} (Theorem~\ref{ec:thm:compat}); \lean{ex341} (Proposition~\ref{ec:prop:ex341}, a finite
  computation).
\item \dirn{PRIMEFOLD/compat/lean/L0Check.lean}: \lean{arith} (Lemma~\ref{ec:lem:arith}); \lean{orbit_dvd}
  (Lemma~\ref{ec:lem:orbit-dvd}); \lean{L0}, \lean{L0_rot} (Theorem~\ref{ec:thm:l0}).
\end{itemize}
Not covered, deliberately: lines 1--1147 of \dirn{LEANEXT/es7corr/ES7Corollaries.lean} (outside the scope of this
section), a verbatim copy of lines 1--1147 of \dirn{LAYERLEAN/LayerLemma.lean}, which belongs to the layer-lemma
sections (we use Theorem~\ref{ll:thm:reduction} and Lemma~\ref{lg:lem:centre} from there); the \texttt{\#print axioms}
commands at the ends of the four other files (\dirn{ES8SYM/lean/C6check.lean}, lines 119--123;
\dirn{ES8SYM/lean/PolarOrient.lean}, lines 127--132; \dirn{PRIMEFOLD/compat/lean/PCompatCore.lean}, lines 285--291;
\dirn{PRIMEFOLD/compat/lean/L0Check.lean}, line 152), which produce the logs cited above and are not results. Out of
scope for this section: the Lean files of \dirn{ES8SLACK} (\dirn{ES8SLACK/review/C6_early_vertex/EarlyVertex.lean},
\dirn{ES8SLACK/review/C2_identity/Identity.lean}, and the two files that according to \dirn{ES8SLACK/README.md} have
not been compiled, \dirn{ES8SLACK/review/C1_certificates/Composition311.lean} and
\dirn{ES8SLACK/review/C3_residual/TemplateRowsC3.lean}); written proofs of the \dirn{ES8SLACK} lemmas are in
\dirn{ES8SLACK/proofs.md}. Fact~\ref{ec:fact:prim}, Corollary~\ref{ec:cor:planar} and the example in
Remark~\ref{ec:rem:compat}(ii) are written steps and not declarations of the five Lean files; only the Fact has Lean
counterparts, the Mathlib lemmas named in its proof.
\endgroup
